\documentclass[11pt,a4paper]{article}
\usepackage[margin=1in]{geometry}
\usepackage{hyperref}
\usepackage{enumitem}
\usepackage{graphicx}
\usepackage{xcolor}
\usepackage{titling}

% -----------------------------------------------------
% Custom formatting commands (simplified, self-contained)
% -----------------------------------------------------

\newcommand{\BvrLabInstructorValue}{Dr. Raghav B. Venkataramaiyer}

\pretitle{\begin{center}\bfseries\huge}
\newcommand{\BvrSubtitle}[1]{%
  \posttitle{%
    \par\end{center}
    \begin{center}\large#1\end{center}
    \vskip 0.5em}%
}

\newcommand{\BvrHint}[1]{%
  {\normalfont\normalsize\itshape\hfill #1}}

% -----------------------------------------------------

\title{StepUp App}

\BvrSubtitle{%
  Project Proposal for UCS503P  \\
  Submitted to: \BvrLabInstructorValue \\ \bfseries
  Thapar Institute of Engineering and Technology%
}

\author{
  Bhavya Jindal, CSED%
  \thanks{Roll No: \texttt{1024030295}}
  \and
  Akshat Batra, CSED%
  \thanks{Roll No: \texttt{1024030292}}
  \and
  Yash, CSED%
  \thanks{Roll No: \texttt{1024030448}}
}

\date{\today}

\begin{document}
\maketitle

\section{Higher-order goal%
  \BvrHint{\ldots secondary framing}}
This project aims to support healthier, more socially connected student
life on campus by making physical activity tracking trustworthy and by
making it easier for students to find reliable sport partners. While
the topic is broader than any single feature, the immediate engineering
objective is to translate \emph{trustworthy activity tracking} and
\emph{effective sport-group formation} into a measurable, deployable
software solution: the StepUp App.

\section{Time-to-value %
\BvrHint{\ldots primary heuristic}}
\begin{itemize}[leftmargin=0.7cm]
    \item We will deliver meaningful increments quickly by prototyping the core user workflow of StepUp and validating it with users within a short iteration window.
    \item We will prioritize fast feedback over perfection by using weekly iterations (and targeting CI-backed automation early) to reduce release friction.
    \item Success is defined by what can be delivered and measured in the first iteration, then improved based on user feedback.
\end{itemize}

\section{Problem Statement}
Campus fitness and sport-partner discovery on Thapar's campus currently
suffer from two distinct gaps:
\begin{itemize}[leftmargin=0.7cm]
    \item \textbf{Unverified activity data:} existing fitness-tracking approaches trust imported step and activity data without checking whether it reflects genuine walking or running. Records may instead come from phone shaking, duplicate syncs, or movement while travelling in a vehicle, which undermines any challenge or leaderboard built on top of that data.
    \item \textbf{Sport-partner discovery as a filtered list, not a match:} students looking for regular sport partners or groups (e.g., badminton, football, running) currently rely on informal channels such as WhatsApp groups, rather than a system that actually matches them on sport, skill, availability, group size, location, and reliability.
\end{itemize}

\textbf{Target users:} Thapar students who want to (a) participate in
campus fitness challenges with data they can trust, and (b) find
regular sport partners or groups for sports such as badminton, football,
and running.

\textbf{Why this matters now (contextual relevance):} Any gamified
fitness or matchmaking feature is only as useful as the trust students
place in it; unverified activity data and ad-hoc partner discovery both
directly undermine that trust, which motivates addressing both problems
as first-class engineering concerns rather than afterthoughts.

\section{Proposed Solution %
\BvrHint{\ldots Software Proposition}}
\subsection{Overview}
Build the \textbf{StepUp App}, consisting of an Android application and
a web dashboard, with the following components:
\begin{itemize}[leftmargin=0.7cm]
    \item \textbf{Android application (student-facing):} activity synchronization via Health Connect, an on-device activity-integrity check, buddy matching requests, and challenge participation.
    \item \textbf{Activity-integrity engine:} distinguishes genuine walking/running from phone shaking, duplicate/repeated sync records, and vehicle movement, running mostly on-device.
    \item \textbf{Sport group matching:} treats sport-buddy formation as a group-formation/matching problem based on sport preference, skill level, availability, group size, coarse location, and reliability, rather than a filtered contact list.
    \item \textbf{Web dashboard:} challenge configuration, pilot monitoring, and aggregate analytics (participation trends, integrity-flag rates, match success rates).
\end{itemize}

\subsection{Core workflow}
\begin{enumerate}[leftmargin=0.7cm]
    \item Student links activity data through Health Connect on the Android app.
    \item The on-device activity-integrity check evaluates incoming activity records (genuine walking/running vs.\ phone shaking, duplicate sync, or vehicle movement) before they count toward a challenge.
    \item Student optionally submits a sport-buddy request specifying sport, skill level, available time slots, preferred group size, and coarse location.
    \item The backend matching engine proposes a group or partner based on these constraints and each student's reliability history.
    \item Admins/organizers use the web dashboard to configure challenges, monitor the pilot, and review aggregate analytics.
\end{enumerate}

\subsection{Operational constraints}
\begin{itemize}[leftmargin=0.7cm]
    \item Activity data will be synchronized via Health Connect (not raw continuous GPS tracking).
    \item The activity-integrity check should run primarily on-device to minimize data transmission and enable near-real-time feedback.
    \item Matching will use only coarse location (e.g., locality/hostel block), not precise GPS coordinates.
    \item The Android application will be built using Kotlin.
\end{itemize}

\section{Solution Approach %
\BvrHint{\ldots Engineering focus}}
This is an engineering course project, so we focus on scalable
implementation, workflow correctness, and rigorous evaluation against
simple baselines, rather than claiming state-of-the-art research
contributions.

\subsection{Core technical approach %
\BvrHint{\ldots non-research}}
\textbf{Activity integrity:} we will build a labelled dataset using
short recorded sessions collected from team members' phones, covering
genuine walking, genuine running, phone shaking/fake steps, duplicate or
repeated sync records, and movement while travelling in a vehicle, using
accelerometer and step data obtained through Health Connect. A
rule-based baseline will use thresholds on step cadence, variance, and
speed. A proposed detector will then be developed and compared against
this baseline; we will not claim any specific accuracy figures until
this evaluation is actually carried out.

\textbf{Sport group matching:} sport-buddy formation will be treated as
a group-formation/matching problem, with inputs including sport
preference, self-reported skill level, available time slots, preferred
group size, coarse location, and reliability based on previous
no-shows. The proposed matching approach will be compared against two
baselines: random matching and first-fit matching.

\subsection{System architecture %
\BvrHint{\ldots solution architecture}}
\begin{itemize}[leftmargin=0.7cm]
    \item \textbf{Android application:} built in Kotlin; handles Health Connect activity synchronization, the on-device activity-integrity check, buddy matching requests, and challenge participation.
    \item \textbf{Web dashboard:} supports challenge configuration, pilot monitoring, and aggregate analytics, including participation trends, integrity-flag rates, and match success rates.
    \item \textbf{Backend:} hosts the matching engine, exposes APIs consumed by both the Android application and the web dashboard, and performs server-side integrity-model comparison and evaluation logging.
    \item \textbf{AI/ML layer:} covers the activity-integrity dataset, the rule-based baseline, the proposed detector, and model evaluation; supports both on-device inference and backend batch evaluation/logging. Specific backend framework and database technology choices are to be finalized during the design phase.
\end{itemize}

\subsection{CI/CD and fast delivery enabler}
CI/CD will be used to:
\begin{itemize}[leftmargin=0.7cm]
    \item Automatically build and run tests on every change.
    \item Produce repeatable deployment artifacts to staging.
    \item Enable safe and frequent releases of incremental features.
\end{itemize}

\section{Evaluation Criterion %
\BvrHint{\ldots measurable and attributable}}
We will define success using measurable metrics tied to the two core
problems: activity integrity and sport group matching.

\subsection{Primary evaluation metric}
\textbf{F1 score of the proposed activity-integrity detector, compared
against the rule-based baseline}, measured on our labelled dataset.
\begin{itemize}[leftmargin=0.7cm]
    \item Target: the proposed detector should achieve a higher F1 score than the rule-based baseline. Specific target values will be set once baseline results are available.
    \item Attribution: computed by comparing detector predictions against ground-truth labels on held-out recorded sessions.
\end{itemize}

\subsection{Secondary evaluation metrics}
\begin{itemize}[leftmargin=0.7cm]
    \item \textbf{Activity integrity:} precision, recall, false-positive rate, on-device latency, and battery cost of running the detector.
    \item \textbf{Matching quality:} match acceptance rate, fairness across users, match stability, and successful completed sessions versus formed matches, each compared against random matching and first-fit matching baselines.
    \item \textbf{Participation:} repeat participation observed in weeks 2--3 of the pilot compared with week 1.
\end{itemize}

\subsection{Pilot validation plan %
\BvrHint{\ldots high-level}}
\begin{itemize}[leftmargin=0.7cm]
    \item Recruit volunteers from our batch/department and hostel/class groups, targeting approximately 40--60 students.
    \item Run the pilot for 2--3 weeks, covering at least one complete challenge cycle.
    \item Track activity-integrity flags, match acceptance/no-shows, and repeat participation throughout the pilot window to evaluate the metrics above.
\end{itemize}

\section{Scalability %
\BvrHint{\ldots with foresight}}
The proposition should scale in both theoretical and practical senses.

\begin{enumerate}[leftmargin=0.7cm]
    \item \textbf{Scalable (theoretically):} The system should support growth in number of students, challenges, and concurrent matching requests by:
    \begin{itemize}[leftmargin=0.7cm]
        \item decoupling the Android application, web dashboard, and backend,
        \item enabling pagination and indexing for common queries (e.g., challenge participation, match history),
        \item using stateless API design for the backend.
    \end{itemize}

    \item \textbf{Operationally deployable (practically):} The system should run on commonly available hosting/deployment infrastructure:
    \begin{itemize}[leftmargin=0.7cm]
        \item managed database,
        \item containerized or platform-hosted backend deployment,
        \item CI/CD pipeline for staging/production.
    \end{itemize}
\end{enumerate}

\section{Engine availability heuristic}
A set of mature tools/services is expected to support this project:
\begin{itemize}[leftmargin=0.7cm]
    \item Health Connect API for activity data synchronization on Android.
    \item Kotlin/Android SDK for the student-facing application.
    \item A backend web framework for REST APIs (specific choice to be finalized).
    \item A managed relational or document database (specific choice to be finalized).
    \item An authentication approach (token-based or session-based, to be finalized).
    \item A test framework for unit/integration tests.
    \item A CI runner and staging deployment target.
\end{itemize}
We will use these components to focus on the activity-integrity and
matching problems, and on measurement, rather than inventing new
infrastructure.

\section{Project scope and deliverables %
\BvrHint{\ldots iteration-friendly}}
\subsection{Initial deliverable %
\BvrHint{\ldots first iteration}}
\begin{itemize}[leftmargin=0.7cm]
    \item Labelled activity-integrity dataset (genuine walking, genuine running, phone shaking, duplicate sync, vehicle movement) and rule-based baseline.
    \item Android application skeleton with Health Connect synchronization and challenge participation flow.
    \item Initial matching formulation with random and first-fit baselines.
    \item Basic logging for activity-integrity and matching metrics.
    \item CI: build + backend unit tests + linting.
    \item CD to staging with a single command or pipeline job.
\end{itemize}

\subsection{Subsequent deliverables}
\begin{itemize}[leftmargin=0.7cm]
    \item Proposed activity-integrity detector, integrated on-device, and its evaluation against the rule-based baseline.
    \item Refined sport-group matching approach, evaluated against random and first-fit baselines.
    \item Web dashboard for challenge configuration, pilot monitoring, and aggregate analytics.
    \item Privacy controls: consent flow, data/match-history deletion, and visibility controls.
    \item Execution of the campus pilot and reporting of results against the evaluation metrics.
\end{itemize}

\section{Team Responsibilities %
\BvrHint{\ldots proposed allocation}}
The following is our proposed division of responsibilities for the
project. This reflects planned ownership going forward and does not
imply that any of this work has already been completed.
\begin{itemize}[leftmargin=0.7cm]
    \item \textbf{Bhavya Jindal:} Backend development, system integration across the Android application, web dashboard, and backend, and overall software engineering coordination.
    \item \textbf{Akshat Batra:} Android application development, including Health Connect integration and integration of the activity-integrity check into the student-facing app.
    \item \textbf{Yash:} AI/ML work for the activity-integrity model, including dataset creation and evaluation, along with support for the web/admin analytics views.
\end{itemize}

\section{Risks and mitigations}
\begin{itemize}[leftmargin=0.7cm]
    \item \textbf{Data privacy / consent:} require explicit consent before collecting activity or location data; do not store continuous GPS trails and use only coarse location for matching; process and discard raw sensor data on-device where possible; provide data deletion, match-history deletion, and visibility controls so students choose what matched buddies can see.
    \item \textbf{User adoption:} keep the on-device integrity check unobtrusive and the matching request flow simple, so participation in challenges and buddy matching requires minimal friction.
    \item \textbf{Evaluation measurement ambiguity:} finalize the labelled-dataset protocol and define activity-integrity and matching metrics before development begins, to avoid ambiguity when comparing against baselines.
    \item \textbf{Operational issues in pilot:} use staging early; ensure CI/CD produces repeatable deployments; test Health Connect synchronization across a range of devices before the pilot begins.
\end{itemize}

\section{Summary}
This project proposes the StepUp App, a campus fitness and sports-partner
platform for Thapar students that addresses two concrete problems:
verifying whether imported activity data is genuine, and treating
sport-buddy formation as a real matching problem rather than a filtered
list. It meets the course criteria by emphasizing time-to-value through
rapid iterations, implementing CI/CD to support frequent safe releases,
and using measurable evaluation criteria for both the activity-integrity
detector (compared against a rule-based baseline) and the matching
approach (compared against random and first-fit baselines), validated
through a campus pilot of approximately 40--60 students.

\end{document}
